\documentclass[11pt]{article}
\usepackage[a4paper,margin=25mm]{geometry}
\usepackage{amsmath,amssymb,amsthm,hyperref}
\newtheorem{theorem}{Theorem}
\newtheorem{lemma}[theorem]{Lemma}
\newtheorem{proposition}[theorem]{Proposition}
\newcommand{\ddc}{\mathrm d'\mathrm d''}
\newcommand{\Int}{\operatorname{Int}}
\newcommand{\ord}{\operatorname{ord}}
\newcommand{\Div}{\operatorname{Div}}
\newcommand{\supp}{\operatorname{supp}}
\title{Ordinary smooth two-sign regularization on a supplied semistable tree}
\author{Unverified research draft for P06.2}
\date{4 October 2026}
\begin{document}
\maketitle

\noindent\textbf{Status and scope.}
This is a conditional construction on an actual supplied smooth affinoid
and actual supplied semistable tree. It does not solve original P06.2 or
the all-curve regularization gate. No current transport, model existence,
normalization, or recognition from a continuous current equation is proved
or assumed. Source extraction and compilation are not mathematical
certification. All arguments below are author drafts.

\begin{theorem}\label{thm:main}
Let $C$ be algebraically closed, complete, and nontrivially valued, in any
characteristic. Let $Z$ be a connected rig-smooth strict affinoid
$C$-analytic curve, with a supplied simply semistable formal model
$\mathfrak Z$, finite metric tree $\Gamma$, actual Shilov boundary $B$,
and retraction $\tau:Z\to\Gamma$. Let $F$ be finite, real, continuous,
affine on the edges, and harmonic on $\Gamma\setminus B$.
Put $h=F\circ\tau$ and $\Omega=\Int(Z)$, the absolute analytic interior.
For each sign $\sigma\in\{1,-1\}$ there are functions
$u_n^\sigma\in C(Z,\mathbb R)$ such that their restrictions to $\Omega$
are ordinary CLD-smooth and psh, $u_{n+1}^\sigma\leq u_n^\sigma$ on $Z$,
and $u_n^\sigma\to\sigma h$ uniformly on $Z$.
Consequently, for every supplied compact strict affinoid
$Y\subset\Omega$, both sequences restrict to decreasing uniformly
convergent sequences on that same $Y$. Every approximant is ordinary
smooth on the actual ordinary open neighborhood $\Omega$ of $Y$.
For every nonnegative compactly supported ordinary smooth scalar test
$\varphi$ on $\Omega$,
\[
 \langle\ddc[u_n^\sigma],\varphi\rangle
 =\int_\Omega u_n^\sigma\ddc\varphi
 =\int_\Omega\varphi\ddc u_n^\sigma\ \geq 0.
\]
No harmonic ordinary-smooth approximants are asserted.
\end{theorem}

\paragraph{Foundations, with their scopes and uses.}
We use the following standard analytic and algebraic foundations.
These are precisely stated inputs; the new constructions and their
applications are proved below. The authoritative companion
\texttt{P06.2-argument.md} includes the selected complete primary proof
units, literal LF ranges, hashes, and the full body of this paper.

\begin{enumerate}
\item A strict affinoid algebra is Noetherian, and a maximal-ideal
residue field is finite over the ground field (Tate affinoid
Noetherianity and Nullstellensatz; \cite[\S6.1.2]{BGR}). A rig-smooth
affinoid over an algebraically closed field is regular. In dimension
one, a connected such algebra is an integral normal ring: distinct
irreducible components of a regular Noetherian ring cannot meet, and
its finitely many components would otherwise give a disconnected
spectrum. For a nonzero element $a$ its zero locus is a finite set of
rigid points. Indeed, in the dimension-one integral ring $A$, $A/(a)$
is Noetherian of dimension zero, hence Artinian with finitely many
maximal ideals; their residue fields are $C$. These are the only
affinoid-algebra facts needed for finiteness of exceptional directions.
For $x\in\mathcal M(A)$, $\mathcal H(x)$ is by definition the completion
of $\operatorname{Frac}(A/\ker x)$; see
\cite[\S2.1.1, immediately before Proposition 2.1.1]{Thu}.

\item For a supplied simply semistable formal curve, its vertices are
the generic points of the reduced special-fibre components, and its
edges are the nodes. A vertex outside the actual analytic boundary has
a proper smooth residual component. For a smooth strict affinoid curve
the analytic boundary is its Shilov boundary, so $\Omega=Z\setminus B$.
The residual component at a type-2 vertex $v$, completed if necessary,
has function field $\widetilde{\mathcal H(v)}$.
The components of $Z\setminus\Gamma$ retract to single points of
$\Gamma$. At a model vertex, the directions are the residual closed
points: nodes give the graph directions; the other points give full
open residue discs. The node tube, with its two endpoints omitted,
is a full ordinary open annulus and retracts to that edge interior.
Sources are \cite[Proposition 2.1.2, Proposition 2.1.12,
\S2.2.2--2.2.3, and the proper-component criterion in the proof of
Proposition 2.2.24]{Thu} and \cite[15.2]{GJR}.
Their hypotheses hold for the actual $\mathfrak Z$; no model is being
manufactured for a different affinoid.

To check the ambient hypothesis of the cited boundary identity
explicitly, the rigid affinoid of $Z$ is normal, connected, and
one-dimensional over complete algebraically closed $C$.
\cite[Theorem 1.1]{VdP}, with precisely those hypotheses, embeds
it as an affinoid domain of a nonsingular complete algebraic curve.
\cite[Theorem 1.6.1 and Lemma 1.6.2(i)]{Ber}, over a nontrivially
valued field and for Hausdorff strict analytic spaces, identifies
this same affinoid with a Berkovich affinoid domain in that curve's
boundaryless analytification. The image is the same affinoid
algebra, not a substituted neighborhood. Thus Thuillier
Propositions 2.1.10 and 2.1.12 apply with that boundaryless ambient
curve: $B=\partial Z$ and its complement is an actual ordinary
open interior. This use of a published algebraization foundation
only checks the boundary identity. It supplies no new formal model,
tree, factorization, or regularization theorem, and no later
argument uses an algebraic presentation of $Z$.

For clarity, the disc/annulus coordinate assertion used here follows
from the given formal charts. At a smooth closed point the chart is
formally etale over a formal line. At a node it is formally etale over
$\operatorname{Spf}(C^\circ\{X,Y\}/(XY-a))$.
The residue field of the selected point is $\widetilde C$. Here is
the actual tube application, including its ordinary domain. The
etale Jacobian-presentation criterion
(\cite[IV, 17.6]{EGA}, applied to the finite-presentation formal
charts) gives, after localization at the selected point, finitely
many equations in auxiliary variables with a square Jacobian
invertible at its residual solution. All localizing denominators
are nonzero at that solution. Choose constant lifts of its residual
coordinates. On any strict closed subdisc $|T|\leq r<1$, or strict
closed subannulus inside $|a|<|T|<1$, the equations evaluated at
these lifts have Banach norm at most some $q<1$. Their Jacobian
has inverse of norm at most one: its residual matrix is invertible,
and the inverse is obtained from its adjugate and the geometric
series for a determinant of the form a unit constant plus an error
of norm less than one. Localization denominators similarly have
norm-one inverses throughout the chosen branch. The base variables
have norm less than one on these strict subdomains; the constant
residual errors have norm less than one. Thus the asserted $q<1$
follows directly from the bounded convergent formal coefficients.

Banach Newton iteration now gives analytic auxiliary coordinates
on that entire subdomain: the first correction has norm at most
$q$, the next error at most $q^2$, and at step $j$ the error is
at most $q^{2^j}$. The coefficients are bounded by one, so the
Taylor remainder has this quadratic bound without dividing by
factorials, in any characteristic. The limit is unique among
solutions having the selected residual coordinates. To check this
last point even for a solution initially farther than $q$ from
the chosen constant lifts, the linear term has the norm of that
distance while the quadratic remainder has strictly smaller norm;
the equation therefore first forces distance at most $q$.
For two such solutions the same invertible linear-term estimate
forces equality. The argument also holds after evaluation in any
completed residue field. Hence the analytic inverse constructed by
Newton is both a left and right inverse on the selected tube,
including nonrigid analytic points, rather than just a bijection
on $C$-points. On overlapping subdomains these inverses agree by
uniqueness. Their ordinary interiors exhaust the disc or annulus,
because the value group of $C$ is dense. This gives respectively
the full ordinary opens $|T|<1$ and $|a|<|T|<1$, with their
entire analytic fibres, as claimed. This proof uses neither an
adic completion at the non-finitely-generated maximal ideal of
$C^\circ$ nor an assertion about the entire etale chart.
The smooth residual point has the parameter $T$ supplied by that
formal chart, normalized so the disc is $|T|<1$ and its added
boundary point is $v$.

\item If $f$ is analytic on the generic fibre of a simply semistable
formal curve and is generically invertible, then its skeletal
$\log|f|$ is integral piecewise affine; on each full retraction fibre
it is an upper bound, with equality if there are no zeros on that
fibre. Its graph Laplacian in the interior is the divisor degree in
that fibre. This is \cite[Proposition 2.2.24]{Thu}. In particular,
for a unit on a standard closed annulus the logarithm is affine on
its skeleton and equals that affine function on every full fibre.
For a connected simply semistable formal curve with exactly one node,
and a unit on its entire generic fibre, the outgoing slope at a
vertex is
\[
                 d_p\log|f|(v)=-\ord_p(\widetilde f),
\]
where $|f(v)|=1$ and $\widetilde f$ is its meromorphic residual
function. This is \cite[Lemma 2.2.25]{Thu}. We apply the lemma only
to actual one-node formal open subcurves on whose generic fibres
the chosen quotient is a unit; that application is constructed below.
The full Proposition 2.2.24 and Lemma 2.2.25 proofs are supplied in
the companion argument, not just their statements.

\item Type-2 cuts in the supplied model are available by admissible
blowups (\cite[Lemma 2.2.22]{Thu}). The construction needed here is
the explicit one in that proof: in a disc direction with formal
parameter $T$, blow up $(T,c)$, $0<|c|<1$; at a node subdivide by
the corresponding $(X,c)$ blowup. The new adjacent node tube is
$|c|<|T|<1$. The strict transform of the original residual component
is unchanged; the new node is its original point $p$. Finitely many
such cuts may be made successively. These are all type-2 cuts.
Any finite extension appearing in that lemma is trivial over $C$.
Retraction compatibility is \cite[Proposition 2.2.26]{Thu}; the
original $\tau$ is retained in the approximation.

\item For a smooth projective connected curve $R$ of genus $g$ over an
algebraically closed field, Riemann--Roch states
$\ell(D)-\ell(K_R-D)=\deg D+1-g$, with $\deg K_R=2g-2$
(\cite[IV, Theorem 1.3 and Corollary 1.4]{Hart}). In particular,
when $\deg D\geq2g$, evaluation
$H^0(R,\mathcal O(D))\to\mathcal O(D)|_p$ is onto for every closed
point $p$. The application and exact-order avoidance are proved below.
No assertion that every degree-zero divisor is principal is used.
The residual component at the interior vertex meets the projective
hypothesis as well as the proper one: an abstract nonsingular curve
is an open subcurve of a nonsingular projective curve
(\cite[I, Theorem 6.9]{Hart}); if the original curve is proper,
its image is also closed in that separated projective curve, and
connectedness makes the image the whole curve. Thus the actual
smooth proper component $R_v$ is itself projective. This check
does not require projectivity of the supplied analytic affinoid.

\item Ordinary smooth functions are locally smooth functions of the
logarithms of a finite tuple of actual analytic units on an ordinary
open neighborhood (\cite[Proposition 15.1.1(b) and Corollary 15.1.2]{CLD}).
Pullback commutes with $\mathrm d',\mathrm d''$, and a positive
top-degree smooth form integrates nonnegatively for the canonical
analytic calibration (\cite[\S7.3.1--7.3.4]{CLD}). On an analytic
interior the calibrated space is harmonious
(\cite[Theorem 15.6.1]{CLD}); the compact-support Stokes and Green
formula, including the case where only the intersection of the
supports is compact, is \cite[Proposition 5.11.1]{CLD}.
These foundations apply to the good, Hausdorff, pure-dimensional,
boundaryless ordinary open $\Omega$ of our affinoid. They imply that
for a smooth $g$ and a compact scalar test $\varphi$ there,
$\int g\ddc\varphi=\int\varphi\ddc g$.
We prove positivity by calculating $\ddc g$ in the actual unit charts,
not by a recognition theorem, weak smoothness, or a curve pullback.
\end{enumerate}

\paragraph{Field and multiplicity checks.}
The residue field $k=\widetilde C$ is algebraically closed: lift a monic
polynomial over $k$ to a monic polynomial over $C^\circ$, factor over
$C$, and use leading-term dominance to see that every root has norm
at most one. Reduction of the factorization proves the assertion.
It is consequently infinite, even if it is the algebraic closure of
a finite field. Taking roots in $C$ makes $|C^\times|$ divisible;
the logarithms of the rational powers of one nontrivial value are
dense in $\mathbb R$. Thus arbitrarily short strict type-2 cuts exist.
The model is simply semistable, so its special fibre is reduced.
The edge weight is $[\kappa(p):k]$ at the corresponding node
(\cite[\S2.2.3]{Thu}, \cite[15.2]{GJR}), which is $1$ here.
The generic integration cycle of $Z$ has length $1$, since $Z$ is
smooth and reduced. These are computed actual multiplicities, not
an imposed normalization. In a non-coordinate unit moment there can
still be higher tropical multiplicities or inseparable degrees;
they are retained in the canonical calibration in the positivity
calculation below. No prime-to-characteristic assumption occurs.

\begin{lemma}[Exact residual movement]\label{lem:rr}
Let $R$ be a smooth projective connected curve of genus $g$ over an
infinite algebraically closed field $k$. Let $p_1,\ldots,p_d$ be distinct
closed points and $q_i\in\mathbb Q$ with $D_0=\sum_i q_i>0$.
For any finite forbidden set $M$ containing those points there is
an integer $N>0$ clearing all denominators and a rational function
$\bar f\in k(R)^\times$ such that
\[
 \ord_{p_i}(\bar f)=-Nq_i,
 \qquad
 \operatorname{div}(\bar f)=A-\sum_i Nq_i[p_i],
\]
where $A$ is effective, supported outside $M$, and $\deg A=ND_0$.
In particular every additional nonzero order is positive.
\end{lemma}
\begin{proof}
Choose a sufficiently large multiple $N$ of a common denominator with
$ND_0\geq\max\{1,2g\}$, and put $E=\sum_iNq_i[p_i]$.
For any $p$, both $K_R-E$ and $K_R-E+p$ have negative degree.
A negative-degree line bundle has no nonzero section, since the
divisor of a nonzero section would be effective of the same degree.
Riemann--Roch therefore gives $\ell(E)-\ell(E-p)=1$.
The kernel of evaluation at $p$ is precisely $H^0(R,\mathcal O(E-p))$,
so the evaluation kernel is a proper linear subspace.
Over an infinite field, a finite union of proper subspaces of a finite
dimensional vector space is not the whole space: take the product of
nonzero linear forms vanishing on those subspaces, and use induction
on the number of variables to see that a nonzero polynomial cannot
vanish at every $k$-point. Choose a section $s$ outside all kernels
for $p\in M$. The space is nonzero by Riemann--Roch and the degree
bound. Regard $s$ as a rational function $\bar f$ in the definition
of $\mathcal O(E)$. Then $A=\operatorname{div}(\bar f)+E$ is effective
and has coefficient zero at every point of $M$. This proves exact
orders, not just lower bounds. Its degree is $\deg E=ND_0$ because
a principal divisor has degree zero. All zeros not prescribed by
negative $q_i$ are precisely the finite effective divisor $A$.
\end{proof}

\begin{lemma}[Actual unit and a common ordinary star]\label{lem:star}
Let $v$ be an interior model vertex of the supplied $Z$, and let
$\bar f\in k(R_v)^\times$ be any nonzero residual rational function.
There is a quotient $f=a/b$ with $a,b\in\mathcal O(Z)$, an actual
ordinary open neighborhood $W$ of $v$ on which $f$ is an analytic
unit, and finitely many exceptional residual directions $E$, with
the following properties. We have $|f(v)|=1$ and reduction $\bar f$.
Every graph direction has a full open annular collar in $W$ on
which $\log|f|$ is its exact affine outgoing germ, of slope
$-\ord_p\bar f$. Every direction in $E$ that is not a graph
direction has such a collar in its full residue disc. In every
other residue disc $f$ is a unit on the entire disc and $|f|=1$.
The set $E$ may contain zero-order directions. In those directions
the collar logarithm is exactly zero, and may be replaced farther
in the disc by that same constant.
\end{lemma}
\begin{proof}
Write $Z=\mathcal M(A)$. At $v$, the residual function field of the
given component is $\widetilde{\mathcal H(v)}$, by the reduced model
and Proposition 2.1.2 of Thuillier. Lift $\bar f$ to
$z\in\mathcal H(v)^\circ$ of norm one. By the definition of
$\mathcal H(v)$, choose $a/b\in\operatorname{Frac}(A)$ with
$|a(v)/b(v)-z|<1$. The quotient has norm one and exactly the same
residue, not an approximate residue. Here $v$ has zero kernel on
$A$: it is not a rigid zero of any nonzero element of the integral
dimension-one ring $A$, as also explained in Thuillier's proof of
Proposition 2.2.24. Thus the fraction genuinely uses global analytic
functions $a,b$, both nonzero at $v$.

Let $P$ be the union of their zero sets. It is finite, consists of
rigid points, and does not contain $v$. On the ordinary open
$Z\setminus P$, $f$ is an analytic unit. Indeed, around a point
where $|a|,|b|>0$, lower bounds on both absolute values give rational
affinoid neighborhoods on which their inverses are analytic, and
the strict lower-bound loci are ordinary open.
The functions $a,b$ can be scaled by the same scalar so that both
have norm one at $v$; this is possible since a model vertex has
$|\mathcal H(v)^\times|=|C^\times|$. Their quotient reduction is
still exactly $\bar f$.

Choose an open metric star $S$ in $\Gamma$ about $v$, containing
no other model vertex or boundary point. Make every graph arm so
short that it avoids $\tau(P)\setminus\{v\}$. This is possible
because that set is finite. The ordinary open $U=\tau^{-1}(S)$
contains the entire residue disc $D_p$ at each non-graph direction
$p$ of $v$. It also contains full edge collars: it is not just a
union of paths in the skeleton.
Let $E$ contain the graph directions, $\supp\operatorname{div}(\bar f)$,
and the directions of all points in $P\cap\tau^{-1}(v)$.
This set is finite. For each non-graph $p\in E$, write $D_p$ as
$|T_p|<1$, as in the formal-chart foundation. Choose
$r_p\in|C^\times|$, $r_p<1$, larger than the norms of all points
of $P\cap D_p$, and let $K_p=\{|T_p|\leq r_p\}$.
It is an actual compact strict closed disc inside $D_p$, hence
closed in the Hausdorff compact $Z$. Set
\[
                     W_0=\tau^{-1}(S)\setminus\bigcup_{p\in E\text{ non-graph}}K_p.
\]
This is an ordinary open neighborhood of $v$, contained in
$\Omega$, disjoint from $P$. Thus $f$ is one actual unit on $W_0$.

We verify all outgoing slopes with their correct formal hypotheses.
On each graph arm choose a strict type-2 cut inside $S$, before
any point of $P$ in that arm. In each exceptional residue disc
use the cut $|T_p|=r_p$, increasing $r_p$ toward $1$ if needed.
The explicit admissible blowups $(T_p,c_p)$ and the node subdivisions
in Thuillier Lemma 2.2.22 create these finitely many adjacent edges.
For a smooth residual point, the charts $T_p=c_pU$ and $T_pV=c_p$
show directly that the strict transform of $R_v$ is isomorphic to
$R_v$, that the new node on it is $p$, and that the tube of that
node is $r_p<|T_p|<1$. At an existing graph node the analogous
subdivision gives the short original edge collar. None of these
node tubes meets $P$, by the cuts just chosen.

In this single refined formal model, fix one adjacent node. Its
special-fibre node point is not the reduction of a point of $P$,
because its entire tube is disjoint from $P$. Remove the reductions
of the finitely many points in $P$ and all other nodes, and retain
a connected formal open neighborhood of the selected node with
exactly this one singular point. Both components through it are
retained at their generic points. The quotient $f$ is a unit on
its entire generic fibre; its reduction at the $v$ component is
the restriction of the same $\bar f$, and the order at $p$ is
unchanged. Lemma 2.2.25 now applies with all its actual hypotheses
and gives $d_p\log|f|(v)=-\ord_p\bar f$.
This use of a formal open is solely a slope computation; it is
not an assertion that its generic fibre is the ordinary star.
The actual ordinary neighborhood is the explicitly open $W_0$.
Repeating the calculation for the finitely many adjacent nodes
does not change $f$ or that neighborhood.

The equality on whole collars, rather than just skeletal traces,
is also exact. Exhaust each selected open annulus by strict closed
subannuli. The quotient is a unit there. Proposition 2.2.24(ii),(iii)
on the standard annulus gives equality on every full fibre and
zero Laplacian at every interior skeletal point. Its piecewise
integer-affine trace is therefore affine. Overlapping subannuli
identify one affine function on the entire connected collar. Its
intercept is $\log|f(v)|=0$ by continuity toward $v$, and its slope
is the just computed $-\ord_p\bar f$. In a graph arm this uses
the original edge coordinate and its original distance from $v$.
For a residue disc it says, with $t_p=-\log|T_p|$,
\[
                    \log|f|=-\ord_p(\bar f)\,t_p
                   \quad (0<t_p<-\log r_p).
\]
The strict subannuli exhaust because $\log|C^\times|$ is dense.

Finally, if $p\notin E$ is a non-graph direction, neither $a$ nor
$b$ vanishes anywhere in $D_p$. Exhaust $D_p$ by nested concentric
strict closed discs. On each such disc both $f$ and $f^{-1}$ have
their maximum at the unique Gauss Shilov point. The two inequalities
force $|f|$ to be constant on that entire closed disc. The nested
constants agree, hence $|f|$ is constant on $D_p$; its value is
one by continuity to $v$. This proves all remaining directional
claims simultaneously. If an exceptional direction has order zero,
the exact collar formula gives $\log|f|=0$ there even if $a$ or
$b$ has zeros deeper in the disc. A zero-order lift exception is
therefore handled without a fictitious all-disc unit assertion.
One may further shorten each of the finitely many cuts to obtain
any smaller common ordinary star. The same $f$ works throughout.
\end{proof}

\begin{lemma}[An analytic cap with exact open overlaps]\label{lem:cap}
On a full residue disc $D=\{|T|<1\}$ attached at $v$, suppose
an actual unit germ on a common ordinary neighborhood has potential
$c-a t$, $a>0$, on the full collar $0<t=-\log|T|<\ell$.
Given any $\epsilon>0$, this potential can be continued as an
ordinary smooth psh function on $D$ which agrees with that potential
on a smaller full collar, is constant on an ordinary inner open
disc, and differs everywhere from $c$ by less than $\epsilon$.
It can be glued to the original unit germ using genuine ordinary
open overlaps; $v$ does not belong to the cap annulus.
\end{lemma}
\begin{proof}
Let $\rho$ be an even nonnegative $C^\infty$ probability density
supported in $(-1,1)$, and set
\[
 \chi_\eta(x)=\int |x-\eta s|\rho(s)\,ds,
 \qquad M_\eta(x,y)=\frac{x+y+\chi_\eta(x-y)}2.
\]
The convolution is smooth and convex, $|\chi_\eta'|\leq1$,
and $\chi_\eta''\geq0$. Jensen's inequality and the triangle
inequality give
$|x|\leq\chi_\eta(x)\leq|x|+\eta$.
Evenness gives zero first moment, so $\chi_\eta(x)=|x|$ if
$|x|\geq\eta$. Its Hessian is
$\chi_\eta''(x-y)(1,-1)^{\mathsf t}(1,-1)/2$ and its first
partials are nonnegative. Thus $M_\eta$ is smooth convex,
equals $\max(x,y)$ when $|x-y|\geq\eta$, and lies between
that maximum and that maximum plus $\eta/2$.

Choose $r>0$ and then $\eta>0$ such that
\[
 0<\eta<ar/2,\qquad r+\eta/a<\ell,
 \qquad ar+\eta/2<\epsilon.
\]
These are strict inequalities. On the full ordinary open annulus
$A=\{e^{-\ell}<|T|<1\}$ put
\[
                 q(t)=M_\eta(c-at,c-ar).
\]
It equals $c-at$ for $0<t<r-\eta/a$ and the \emph{new}
constant $c-ar$ for $t>r+\eta/a$.
Take $s$ strictly between $r+\eta/a$ and $\ell$, and continue
by $c-ar$ on the ordinary open inner disc $E=\{|T|<e^{-s}\}$.
The two opens cover $D$ and agree exactly on the full open annulus
$s<t<\ell$. To glue at $v$, shorten the star cut to any
$t_0\in(0,r-\eta/a)$. On its overlap with $A$ there is the
entire annulus $0<t<t_0$, with identical values. No partition
gluing is involved. The omitted interface Gauss points are
included in one or the other of these open sets.

The coordinate $T$ is a unit on $A$, and $t=-\log|T|$ is
an actual ordinary moment coordinate. The function $q$ is
smooth convex in $t$, so on $A$
\[
                    \ddc q=q''(t)\,\mathrm d't\wedge\mathrm d''t,
                    \qquad q''(t)\geq0.
\]
On $E$ the Hessian is zero. Hence these are ordinary smooth
functions with positive Hessian form. Their current positivity
against arbitrary compact tests is proved below using the actual
analytic calibration. The bound from the regularized maximum is
$|q-c|\leq ar+\eta/2<\epsilon$, and the constant continuation
has the same bound. In particular we never restore the old
constant $c$. No annular chart containing the positive-genus
vertex was claimed; it is covered by the unchanged unit germ.
\end{proof}

\begin{proposition}[Uniform ordinary approximation on the supplied domain]\label{prop:raw}
Under the geometric hypotheses of Theorem \ref{thm:main}, let $H$
be a finite real continuous edge-affine function on $\Gamma$ with
nonnegative outgoing-slope sum at every model vertex outside $B$.
For every $\epsilon>0$ there is $g\in C(Z,\mathbb R)$, ordinary
CLD-smooth and psh on $\Omega$, with
$\|g-H\circ\tau\|_{C(Z)}<\epsilon$.
\end{proposition}
\begin{proof}
If $\Gamma$ is a point, $H\circ\tau$ is constant and already
works. Otherwise let $V$ be the original model vertices and
$I=V\setminus B$. We use no type-3 model vertices. Redundant
metric subdivisions at which $H$ is harmonic may be removed;
in Theorem \ref{thm:main} the slopes on their two arms are equal.
Choose a root $b\in B$, which is nonempty for a nonempty
affinoid. Set $\mu=1$ on $I$, $\mu(b)=-|I|$, and zero on
$B\setminus\{b\}$. Root the tree at $b$. For each child $x$
let $m_x$ be the sum of $\mu$ over the child subtree; prescribe
slope $m_x$ from $x$ to its parent and integrate with root value
zero to get $\psi$. At $x\ne b$ the outgoing sum is
$m_x-\sum_{y\text{ child of }x}m_y=\mu(x)$; at $b$ the same
identity follows from $\sum\mu=0$. Thus $\Delta\psi=1$ on $I$.

For $\delta>0$ small, $H+\delta\psi$ is within $\epsilon/4$
of $H$ and has outgoing-sum margin at least $\delta$ on $I$.
Round its oriented slopes to rationals with error less than
$\eta_0$, and integrate them from the root's original value to
get $G$. Path uniqueness on the tree gives consistency, and if
$L$ is total length and $d$ maximum valency,
\[
 \|G-(H+\delta\psi)\|_\infty\leq L\eta_0,
 \qquad |\Delta G-\Delta(H+\delta\psi)|\leq d\eta_0.
\]
Choose $L\eta_0<\epsilon/4$ and, if $I$ is nonempty,
$d\eta_0<\delta/2$. Then
$\|G-H\|_\infty<\epsilon/2$ and $\Delta G(v)>0$ on $I$.
When $I$ is empty the outgoing-sum condition is empty.
This elementary tree step is reproved here for use, not presented
as a new analytic result or as a fixed-cycle rounding procedure.

Fix $v\in I$. Its graph directions correspond to the distinct
node points $p_i$ of the smooth proper residual component $R_v$;
there are no self-nodes since the model is simply semistable.
Let $q_i$ be the outgoing slopes of $G$. Their sum is positive.
Apply Lemma \ref{lem:rr}, forbidding every graph direction.
It produces $N_v$ and
$\operatorname{div}(\bar f_v)=A_v-\sum_i N_vq_i[p_i]$,
where $A_v=\sum_{p\in S_v}m_p[p]$ is effective and all its
points are fresh non-graph smooth points. Every such direction
is a full available residue disc of this actual $Z$.
Apply Lemma \ref{lem:star} to obtain one actual unit $f_v$.
The ordinary smooth potential
\[
                   H_v=G(v)+N_v^{-1}\log|f_v|
\]
has exactly the graph slopes $q_i$, exactly the fresh slopes
$-m_p/N_v<0$, and zero slopes in all other directions, including
finite additional lift exceptions. It is an affine function of
an actual unit moment on the entire common star, so its Hessian
form is identically zero there. It is not being inferred smooth
from a weak harmonic tropicalization or from a germ exact sequence.

Choose the metric graph stars at the finitely many $v\in I$ so
short that they are pairwise disjoint and miss $B$. Lemma
\ref{lem:star} allows further shrinking without changing $f_v$.
On every graph collar $H_v=G\circ\tau$ exactly on the whole
annular fibre: intercept and slope agree. In zero-order exceptional
residue directions glue $H_v$ to the constant $G(v)$ on the
whole disc, using the exact constant collar. On every other
non-fresh direction the full-disc potential is already $G(v)$.
For each fresh $p\in S_v$, apply Lemma \ref{lem:cap} with
$c=G(v)$, $a=m_p/N_v$, and error $\epsilon/2$. There are only
finitely many such directions. Make the finitely many star cuts
shorter than their exact affine overlap lengths, while leaving
each cap annulus long enough for its transition and constant
overlap. The common star remains ordinary open because it is
an inverse image of an open graph star minus finitely many
compact closed inner discs. The unchanged central unit germ,
caps, and constant disc continuations therefore agree on all
ordinary open overlaps.

On the full annulus attached to each original open edge put the
original affine value of $G\circ\tau$ in its annulus coordinate.
This is ordinary smooth with zero Hessian. On every residue
disc at a model vertex, use the constant or capped function
just specified; at a boundary vertex use the constant $G(b)$
on its residue discs and $G\circ\tau$ on a small graph star
to check continuity. No smoothness at a boundary vertex is
required. Graph collars at interior vertices agree with the
edge formula on full ordinary open annuli. Residue discs
attached at points in an edge interior are already included
in that edge's full annulus; their values are the same affine
edge value. This accounts for all analytic points, including
the type-3 and type-4 points in the annuli and discs.

The preceding sets give an actual ordinary open cover of
$\Omega$: central stars at interior model vertices, full
edge annuli, and full residue discs (a fresh disc itself is
covered by its cap annulus and inner disc). Exact agreement
on every overlap defines a single ordinary smooth function
there. It extends continuously to $Z$, with $g(b)=G(b)$:
there are finitely many interior stars, they avoid $B$, and
near a boundary vertex the function is just $G\circ\tau$.
The replacement regions are disjoint apart from their exact
overlaps. Outside a fresh residue disc the value is exactly
$G\circ\tau$; inside it the latter is the constant $G(v)$
and the error is less than $\epsilon/2$. Therefore
\[
 \|g-G\circ\tau\|_{C(Z)}<\epsilon/2,
 \qquad \|g-H\circ\tau\|_{C(Z)}<\epsilon.
\]

Here is the positivity check with unrestricted ordinary tests.
In a central star, $g=G(v)+N_v^{-1}\log|f_v|$, so
$\ddc g=0$ as a smooth form. In an edge annulus it is affine
in $-\log|T|$, so again the form is zero. In constant regions
it is zero, and in cap charts it is
$q''(t)\mathrm d't\wedge\mathrm d''t$ with $q''\geq0$.
There are no interface distributions, since these smooth
formulas agree on ordinary open overlaps. To see the sign
in an arbitrary common unit chart for a test and the cap,
on each one-dimensional calibrated cell let $\xi$ be a
primitive tangent coordinate, $m_\alpha>0$ its actual
canonical multiplicity, and $t=A_\alpha\xi+c_\alpha$.
Its contribution is
\[
            m_\alpha\varphi\,q''(t)A_\alpha^2\,d\xi\geq0.
\]
Cells with $A_\alpha=0$ contribute zero, including the
residual directions collapsed by this coordinate. We neither
discard those directions nor replace $m_\alpha$ by one.
The coordinate annulus itself has calibration one, but
refined moment charts can have other multiplicities.
Equivalently, $\ddc g$ is a positive ordinary top-degree
form everywhere, and canonical integration preserves its
positivity (CLD \S7.3.4). For every compact nonnegative
ordinary scalar test $\varphi$ on $\Omega$, Green's formula
with $g$ and $\varphi$ applies because their support
intersection is compact and $\Omega$ is boundaryless:
\[
        \langle\ddc[g],\varphi\rangle
        =\int_\Omega g\ddc\varphi
        =\int_\Omega\varphi\ddc g\geq0.
\]
This proves ordinary current positivity for every test, and
thus psh of the smooth approximant in the ordinary CLD
sense. Only tests, not psh potentials, may be localized by
partitions; the construction did not use partition gluing.
\end{proof}

\begin{proof}[Proof of Theorem \ref{thm:main}]
Both $F$ and $-F$ meet Proposition \ref{prop:raw}'s outgoing-sum
hypothesis with sum zero. For each sign and $b_n=4^{-n}$,
choose $g_n^\sigma$ on the same entire $Z$ with
$\|g_n^\sigma-\sigma h\|_{C(Z)}<b_n$.
Set
\[
 c_n=\sum_{j\geq n}(b_j+b_{j+1})=\tfrac53 b_n,
 \qquad u_n^\sigma=g_n^\sigma+c_n.
\]
Then $c_n-c_{n+1}=b_n+b_{n+1}$, so
$u_n^\sigma-u_{n+1}^\sigma\geq0$ on $Z$.
Furthermore
\[
             0\leq u_n^\sigma-\sigma h\leq\tfrac83 b_n
\]
on $Z$. Thus each sequence decreases uniformly from above.
Constants preserve ordinary smoothness and current positivity.
No domain shrinks with $n$ or with the sign. Restricting to
the supplied $Y\subset\Omega$ proves the required assertion
on precisely that $Y$, with the ordinary neighborhood
$\Omega$ available for every approximant. The functions are
continuous up to $B$, but their boundary values need not
equal $\sigma F|_B$ at any finite stage.
\end{proof}

\paragraph{Comparison with original P06.2.}
The input here already provides smooth reduction, an actual
semistable model with tree skeleton, and harmonic skeletal
factorization over an algebraically closed ground field.
The theorem establishes the ordinary two-sign regularization
construction under those inputs, including positive residual
genus, real slopes, all full annular fibres, and finite lift
exceptions. It does not derive any of those inputs from an
arbitrary continuous zero-current function. It does not prove
the all-curve gate, scalar-current ascent, or the unchanged
good separated strict boundaryless arbitrary-dimensional
recognition atom. No strict analytic counterexample is given.
The catalogue's claimed book regularization result was read
as fallible mathematical source material, not used as a theorem
grant. The tree and cap calculations already present in the
predecessor are reproved as supporting steps, not counted as
new analytic results. The substantive proposed advance is
the explicit residual lift, finite exceptional-direction
control, and the resulting actual whole-overlap analytic
construction. P05.2 is untouched. Candidate remains false.

\begin{thebibliography}{9}
\bibitem{BGR} S. Bosch, U. G\"untzer, R. Remmert,
\emph{Non-Archimedean Analysis}, Springer, 1984, \S6.1.2
(affinoid Noetherianity and Nullstellensatz); regular affinoid
algebras and the Jacobian criterion in \S7.3.
\bibitem{EGA} A. Grothendieck and J. Dieudonn\'e,
\emph{\'El\'ements de g\'eom\'etrie alg\'ebrique IV},
Publ. Math. IH\'ES 32 (1967), \S17.1 and \S17.6
(formal etaleness and etale morphisms).
\bibitem{Hart} R. Hartshorne, \emph{Algebraic Geometry},
Springer, 1977, Chapter IV, Theorem 1.3 and Corollary 1.4.
Chapter I, Theorem 6.9, for projective completion of a nonsingular curve.
\bibitem{Thu} A. Thuillier, \emph{Th\'eorie du potentiel sur
les courbes en g\'eom\'etrie analytique non archim\'edienne.
Applications \`a la th\'eorie d'Arakelov}, thesis, 2005,
HAL tel-00010990v1. Exact installed version, not a later
transcription.
\bibitem{CLD} A. Chambert-Loir and A. Ducros,
\emph{Formes diff\'erentielles r\'eelles et courants sur les
espaces de Berkovich}, arXiv:1204.6277v2, 25 July 2025.
The cited numbering is that of the installed v2 PDF, SHA-256
\texttt{5690f999523017d184ab751bb55857ab1637b197d1186ad4fd2edba7e1bda719}.
\bibitem{GJR} W. Gubler, P. Jell, J. Rabinoff,
\emph{Forms on Berkovich spaces based on harmonic tropicalizations},
arXiv:2111.05741v3, exact installed version, \S15.2.
\bibitem{VdP} M. van der Put, \emph{The class group of a
one-dimensional affinoid space}, Ann. Inst. Fourier 30(4) (1980),
155--164, Theorem 1.1. Its ground field is complete and
algebraically closed.
\bibitem{Ber} V. Berkovich, \emph{\'Etale cohomology for
non-Archimedean analytic spaces}, Publ. Math. IH\'ES 78 (1993),
5--161, Theorem 1.6.1 and Lemma 1.6.2,
DOI 10.1007/BF02712916.
\end{thebibliography}
\end{document}
